\begin{resumo}
O diagnóstico de sobreajuste (\emph{overfitting}) em redes neurais profundas
depende, na prática corrente, de um conjunto de validação rotulado,
representativo e suficientemente grande --- condição frequentemente inviável em
imagens médicas, nas quais a anotação é onerosa e exige especialista. Esta
dissertação investiga se medidas de complexidade estatística calculadas
exclusivamente sobre os pesos da rede, sem acesso a rótulo algum, constituem
marcadores do início do sobreajuste. São consideradas a complexidade
Lopez-Ruiz--Mancini--Calbet (LMC), produto da entropia de Shannon pela
distância estatística ao equilíbrio, a entropia amostral (SampEn) e sua variante
bidimensional (SampEn2D), calculadas a cada época sobre os pesos da camada densa
de redes treinadas no conjunto MedNIST. A memorização é induzida por corrupção
controlada de rótulos, preservando o volume de dados. Com oito execuções e
controles estritamente pareados, a SampEn2D inflete
$7{,}5 \pm 0{,}9$ épocas antes da degradação da acurácia de validação, com
razão de amplitude de $5{,}7$ vezes entre execuções com e sem memorização
($p = 0{,}00016$), resultado replicado na arquitetura DenseNet-121
($4{,}3$ vezes; $p = 0{,}0016$). A complexidade LMC, contudo, não discrimina
memorização na camada densa ($1{,}1$ vez; $p = 0{,}19$). Conclui-se que a
capacidade de marcar o sobreajuste depende da medida empregada.

\textbf{Palavras-chave}: Complexidade estatística; Entropia amostral;
Sobreajuste; Aprendizagem profunda; Imagens médicas; Memorização;
Complexidade LMC.
\end{resumo}
